%============================================================
% CH1 · PATTERN 5 — EXERCISES, HINTS, ANSWERS
%============================================================
\subsection{Pattern 5: Exercises and Going Further}
\label{ssec:ch1-p5-ex}

\begin{patternbox}{Pattern 5 Summary Card}
\begin{tabularx}{\linewidth}{@{}>{\bfseries}l X@{}}
Recognise & small value range (digits, letters, $|a_i| \le 10^3$, $a_i \le n$); the answer depends on values only through their counts \\
Core loop & \texttt{cnt[v - MIN]++} for every element; then scan the $V$ cells \\
Tools & prefix counts $\texttt{less}[v]$ (ranks, counting sort); frequency of frequencies; pairs $\sum_v \texttt{cnt}[v]\cdot\texttt{cnt}[v+k]$; capping values at $n$ \\
Cost & $\bigO{n + V}$ time, $\bigO{V}$ memory; hash map or sorting when $V$ is huge \\
Top traps & $\text{MAX}$ cells instead of $\text{MAX} + 1$; negative values without offset; zero counts treated as occurrences; clearing a big array per test case \\
\end{tabularx}
\end{patternbox}

\exlegend
\begin{enumerate}[leftmargin=*, label=\textbf{E5.\arabic*}, itemsep=6pt]
    \item \exeasy\ \pconcept \exprob{Find Lucky Integer in an Array}{LeetCode 1394}{https://leetcode.com/problems/find-lucky-integer-in-an-array/}\ \lc{Easy}. A value is lucky if it occurs exactly as many times as its value. Return the largest lucky value, or $-1$ ($n \le 500$, $1 \le a_i \le 500$).
    \item \exeasy\ \pconcept \exprob{Count Number of Pairs With Absolute Difference K}{LeetCode 2006}{https://leetcode.com/problems/count-number-of-pairs-with-absolute-difference-k/}\ \lc{Easy}. Count the pairs $i < j$ with $|a_i - a_j| = k$ ($n \le 200$, $1 \le a_i \le 100$, $1 \le k \le 99$). Solve it in $\bigO{n + V}$.
    \item \exeasy\ \precent \cf{800}\ \exprob{Problem Generator}{Codeforces 1980A}{https://codeforces.com/problemset/problem/1980/A}. Vlad wants to hold $m \le 5$ rounds; each round needs one problem of each difficulty \texttt{A}--\texttt{G}. He has $n \le 50$ problems with given difficulties. How many new problems must he create at least?
    \item \exeasy\ \ptricky \exprob{Split the Array}{LeetCode 3046}{https://leetcode.com/problems/split-the-array/}\ \lc{Easy}. Can an array of even length $n \le 100$ ($1 \le a_i \le 100$) be split into two halves of size $n/2$, each consisting of distinct elements?
    \item \exmedium\ \pclassic \exprob{H-Index}{LeetCode 274}{https://leetcode.com/problems/h-index/}\ \lc{Medium}. The h-index of a researcher is the largest $h$ such that at least $h$ papers have at least $h$ citations each. Compute it for $n \le 5000$ papers with up to 1000 citations each, in $\bigO{n}$.
    \item \exmedium\ \ptricky \cf{1100}\ \exprob{Two Teams Composing}{Codeforces 1335C}{https://codeforces.com/problemset/problem/1335/C}. From $n$ students with skills $a_i \le n$, form two disjoint teams of equal size $x$: in the first all skills are distinct, in the second all skills are equal. Maximise $x$.
\end{enumerate}

\subsubsection*{Hints}
\begin{enumerate}[leftmargin=*, label=\textbf{E5.\arabic*}]
    \item One counting array; scan it from the top value down.
    \item For a fixed value $v$, how many pairs use $v$ as the smaller element? Pairs of equal values never qualify since $k \ge 1$.
    \item Seven counters. For each difficulty, how many are missing?
    \item What goes wrong if some value appears three times? And if every value appears at most twice?
    \item A paper with more than $n$ citations is no better than one with exactly $n$. Count papers by citation number, capped at $n$, then sweep $h$ downwards.
    \item Two numbers describe everything: the number of distinct skills $d$ and the largest count $m$. Does the most frequent skill also appear in the distinct team?
\end{enumerate}

\subsubsection*{Answers}

\begin{solutionbox}[E5.1 Find Lucky Integer in an Array]
Count the values (\texttt{cnt} of size 501) and scan $v = 500, 499, \dots, 1$; the first $v$ with $\texttt{cnt}[v] = v$ is the answer. Samples: $[2,2,3,4] \to 2$; $[1,2,2,3,3,3] \to 3$ (1, 2 and 3 are all lucky, 3 is the largest); $[2,2,2,3,3] \to -1$. $\bigO{n + 500}$.
\end{solutionbox}

\begin{solutionbox}[E5.2 Count Number of Pairs With Absolute Difference K]
Every qualifying pair consists of one copy of some $v$ and one copy of $v + k$; the copies can be chosen independently, so the number of such pairs is $\texttt{cnt}[v] \cdot \texttt{cnt}[v+k]$. Each pair is counted once (at its smaller value).
\begin{lstlisting}[style=cpp]
#include <bits/stdc++.h>
using namespace std;

class Solution {
public:
    int countKDifference(vector<int>& nums, int k) {
        vector<int> cnt(101, 0);
        for (int x : nums) cnt[x]++;
        int pairs = 0;
        for (int v = 1; v + k <= 100; v++) pairs += cnt[v] * cnt[v + k];   // smaller value v
        return pairs;
    }
};

int main() {
    vector<int> a = {1, 2, 2, 1}, b = {3, 2, 1, 5, 4};
    cout << Solution().countKDifference(a, 1) << " "
         << Solution().countKDifference(b, 2) << "\n";   // 4 3
}
\end{lstlisting}
\begin{center}
\begin{tikzpicture}[x=1cm, y=1cm]
    \dpindex{vv}{0.65}{value}{1}{5}
    \dparray{c}{0}{\texttt{cnt}}{1/compcell, 1/compcell, 1/compcell, 1/compcell, 1/compcell}
    \foreach \x/\y in {0/2, 1/3, 2/4}{ \draw[dep] (c\x.south) to[bend right=35] (c\y.south); }
    \node[note, anchor=west, text width=4.6cm, align=left] at (5.0,0) {$[3,2,1,5,4]$, $k = 2$: pairs $(1,3), (2,4), (3,5)$, each $1 \cdot 1$ $\Rightarrow 3$};
\end{tikzpicture}
\end{center}
Samples $\to 4, 0, 3$; matches the $\bigO{n^2}$ pair loop on random arrays.
\end{solutionbox}

\begin{solutionbox}[E5.3 Problem Generator]
Count the problems of each difficulty (\texttt{cnt[c - 'A']}, 7 cells). Difficulty $c$ needs $m$ problems, so $\max(0, m - \texttt{cnt}[c])$ must be created; extra problems of one difficulty cannot replace another. Answer: $\sum_c \max(0, m - \texttt{cnt}[c])$.
\begin{lstlisting}[style=cpp]
#include <bits/stdc++.h>
using namespace std;

int main() {
    int t;
    cin >> t;
    while (t--) {
        int n, m;
        string s;
        cin >> n >> m >> s;
        int cnt[7] = {0};                                   // reset per test: only 7 cells
        for (char ch : s) cnt[ch - 'A']++;
        int need = 0;
        for (int c = 0; c < 7; c++) need += max(0, m - cnt[c]);
        cout << need << "\n";
    }
}
\end{lstlisting}
First sample test \texttt{BGECDCBDED}, $m = 1$: counts \texttt{A}..\texttt{G} $= 0, 2, 2, 3, 2, 0, 1$, so \texttt{A} and \texttt{F} are missing $\Rightarrow 2$; with $m = 2$: $2 + 0 + 0 + 0 + 0 + 2 + 1 = 5$; third test $\to 1$.
\end{solutionbox}

\begin{solutionbox}[E5.4 Split the Array]
\textbf{The trap:} it looks like a search over splits, but only the counts matter. A value that occurs three or more times cannot be split into two halves without a repetition (pigeonhole) $\Rightarrow$ \texttt{false}. If every count is at most 2, put one copy of each doubled value into each half; the values occurring once number $n - 2\cdot(\#\text{doubled})$, an even number, so they can be distributed half and half. Answer: \texttt{true} iff $\max_v \texttt{cnt}[v] \le 2$.
\begin{lstlisting}[style=cpp]
#include <bits/stdc++.h>
using namespace std;

class Solution {
public:
    bool isPossibleToSplit(vector<int>& nums) {
        vector<int> cnt(101, 0);
        for (int x : nums)
            if (++cnt[x] > 2) return false;    // a third copy cannot be placed
        return true;
    }
};

int main() {
    vector<int> a = {1, 1, 2, 2, 3, 4}, b = {1, 1, 1, 1};
    cout << Solution().isPossibleToSplit(a) << Solution().isPossibleToSplit(b) << "\n";   // 10
}
\end{lstlisting}
Samples $\to$ true, false.
\end{solutionbox}

\begin{solutionbox}[E5.5 H-Index]
The answer is at most $n$, so citations above $n$ can be capped at $n$ without changing it. Count $b[c]$ = number of papers with $\min(\text{citations}, n) = c$, then sweep $h = n, n-1, \dots$ keeping $\text{atLeast}(h) = \sum_{c \ge h} b[c]$ (a suffix sum of the counts); the first $h$ with $\text{atLeast}(h) \ge h$ is the h-index.
\begin{lstlisting}[style=cpp]
#include <bits/stdc++.h>
using namespace std;

class Solution {
public:
    int hIndex(vector<int>& citations) {
        int n = citations.size();
        vector<int> b(n + 1, 0);
        for (int c : citations) b[min(c, n)]++;      // cap: more than n is as good as n
        int atLeast = 0;                             // papers with >= h citations
        for (int h = n; h >= 0; h--) {
            atLeast += b[h];
            if (atLeast >= h) return h;
        }
        return 0;
    }
};

int main() {
    vector<int> a = {3, 0, 6, 1, 5}, c = {1, 3, 1};
    cout << Solution().hIndex(a) << " " << Solution().hIndex(c) << "\n";   // 3 1
}
\end{lstlisting}
\begin{center}
\begin{tikzpicture}[x=1cm, y=1cm]
    \dpindex{hh}{0.65}{$c$ (capped)}{0}{5}
    \dparray{b}{0}{$b$}{1/compcell, 1/compcell, 0/compcell, 1/compcell, 0/compcell, 2/compcell}
    \dparray{s}{-1.1}{atLeast}{5/basecell, 4/basecell, 3/basecell, 3/anscell, 2/basecell, 2/basecell}
    \node[note, anchor=west, text width=4.2cm, align=left] at (5.8,-0.5) {$[3,0,6,1,5]$, $n = 5$: the 6 is capped to 5. From the right: $h = 5$: $2 < 5$; $h = 4$: $2 < 4$; $h = 3$: $3 \ge 3$ $\Rightarrow 3$};
\end{tikzpicture}
\end{center}
$\bigO{n}$ time; sorting gives $\bigO{n \log n}$. Checked against trying every $h$ on random inputs.
\end{solutionbox}

\begin{solutionbox}[E5.6 Two Teams Composing]
Let $d$ = number of distinct skills and $m$ = the largest count, of skill $s$. Use skill $s$ for the uniform team (any other choice is no better). Two cases: skill $s$ is not in the distinct team --- then $x \le d - 1$ and $x \le m$; or it is --- then $x \le d$ and only $m - 1$ copies remain, $x \le m - 1$. Answer $\max(\min(d-1, m),\ \min(d, m-1))$. \textbf{The trap} is $\min(d, m)$, which forgets that one copy of $s$ may be needed in both teams: for $[1,1,1,3]$ ($d = 2$, $m = 3$) it gives 2, which is right, but for $[4,2,4,1,4,3,4]$ ($d = 4$, $m = 4$) it gives 4, while the answer is 3.
\begin{lstlisting}[style=cpp]
#include <bits/stdc++.h>
using namespace std;

int main() {
    int t;
    cin >> t;
    while (t--) {
        int n;
        cin >> n;
        vector<int> cnt(n + 1, 0);                  // skills are <= n: size by the test
        int d = 0, m = 0;
        for (int i = 0; i < n; i++) {
            int a;
            cin >> a;
            if (cnt[a]++ == 0) d++;                 // a new distinct skill
            m = max(m, cnt[a]);
        }
        cout << max(min(d - 1, m), min(d, m - 1)) << "\n";
    }
}
\end{lstlisting}
Samples: $(d, m) = (4,4) \to 3$, $(5,1) \to 1$, $(1,1) \to 0$, $(2,3) \to 2$. The counting array is sized $n + 1$ per test, so many small tests stay cheap (Tool~5). Checked against trying every team size and every choice of the uniform skill for $n \le 8$.
\end{solutionbox}
